\section{Classes}

\begin{frame}{What are classes?}
\pause
\begin{block}{Class (definition)}
    A class is a blueprint for a type of object, defining it's \textbf{data attributes} and \textbf{methods}.
\end{block}
\pause
\begin{alertblock}{Terminology}
    We call an object from a certain class, an \textbf{instance} of that class. E.g. "OOP" is an instance of the class str (string).
\end{alertblock}
\end{frame}

\begin{frame}{Why do we use classes?}
    \pause
    \begin{itemize}
        \pause
        \item Python only gives us a limited amount of types, integers, floats, strings, etcetera. 
        \pause
        \item With classes we can extend this and model anything we want!
        \pause
        \begin{itemize}
            \item Like students, bank accounts, courses, shopping cart.
        \end{itemize}
        \pause
        \item Encapsulation (later)
        \pause
        \item Seperation of Concerns (later)
    \end{itemize}
\end{frame}

\begin{frame}{What do classes consist of?}
    \begin{itemize}
        \pause
        \item Methods (behaviour)
        \pause
        \item Data Attributes (properties/fields)
    \end{itemize}
\end{frame}

\begin{frame}{What do classes consist of?}
    \pause
    Data Attributes:
    \begin{itemize}
        \pause
        \item Variables that belong to a single object
        \pause
        \item For example a person can have the following data variables (i.e. data attributes) which we would like to store:
        \pause
        \begin{itemize}
            \item first name (string)
            \item last name (string)
            \item age (int)
        \end{itemize}
        \pause
        \item These variables can be of any type! (Strings, integers, lists, and also objects of other classes)
    \end{itemize}
\end{frame}

\begin{frame}{Methods}
\begin{itemize}
    \item You have already seen a lot of methods
    \pause
    \item For example, for a person we could have:
    \begin{itemize}
        \pause
        \item Speak
        \item Dance
        \item Study
    \end{itemize}
    \pause
    \item Every class needs one special method:
    \begin{itemize}
        \pause
        \item \textbf{The constructor}
    \end{itemize}
\end{itemize}
\end{frame}

\begin{frame}{What is a constructor?}
\pause
\begin{block}{Constructor}
In Python, a \textbf{constructor} is a special method called when an object is created. It initializes all data attributes.
\end{block}
\end{frame}

\begin{frame}[fragile]{The constructor}
To create objects of a certain class:
\pause
\begin{minted}{python}
my_string = "CognAC"  # Class str
my_int = 18  # Class int
my_list = [18, 2003]  # Class list
\end{minted}
\pause
\textbf{But what if we have a new type of object of a new class, for example a Fraction class}
\pause
\textit{We use the constructor!}
\end{frame}

\begin{frame}[fragile]{The constructor}
\begin{minted}{python}
# We import the class from the module 'fractions'
# Classes always start with a capital lette
from fractions import Fraction
\end{minted}{python}
\pause
\begin{minted}{python}
# We call the constructor by 'calling the class',
# this creates a new object for us!
# The constructor takes two arguments:
# numerator and denominator (in this case 8 and 10)
my_fraction = Fraction(8, 10)  # Fraction of 8/10th
\end{minted}
\end{frame}

\begin{frame}{Making our own class!}
\begin{itemize}
    \pause
    \item We now have all ingredients to create our own classes!
    \pause
    \item Let's say we want to model a person in python
    \pause
    \item We want it to have a \textit{first name, last name, age} (data attributes)
    \pause
    \item And we want a person to 'speak' i.e. print some text (methods)
    \pause
    \item We create our own class \textbf{Person}
\end{itemize}
\end{frame}

\begin{frame}[fragile]{Creating our Person class}
\begin{minted}[fontsize=\footnotesize, breaklines]{python}
class Person():
\end{minted}
\pause
\begin{minted}[fontsize=\footnotesize, breaklines]{python}
    
    # The constructor
    def __init__(self, first_name, last_name, age):
        self.first_name = first_name  # Data attribute
        self.last_name = last_name  # Data attribute
        self.age = age  # Data attribute
\end{minted}
\pause
\begin{minted}[fontsize=\footnotesize, breaklines]{python}

    # Method
    def speak(self):
        print(f"Hey, my name is {self.first_name}")
\end{minted}
\end{frame}

\begin{frame}{The 'self' parameter}
\begin{itemize}
    \pause
    \item The first parameter in a constructor definition is always named \textbf{self}. 
    \pause
    \item This refers to the object itself, and is necessary for declaring any data attributes attached to the object.
    \pause
    \item It is automatically filled in by the interpreter with the object it is called on.
\end{itemize}
\end{frame}

\begin{frame}[fragile]{Creating our Person class}
\begin{minted}[fontsize=\footnotesize, breaklines]{python}
class Person():
    
    # The constructor
    def __init__(self, first_name, last_name, age):
        self.first_name = first_name  # Data attribute
        self.last_name = last_name  # Data attribute
        self.age = age  # Data attribute

    # Method
    def speak(self):
        print(f"Hey, my name is {self.first_name}")
\end{minted}
\end{frame}

\begin{frame}[fragile]{Data Attributes}
Accessing a data attribute:
\pause
\begin{minted}{python}
person = Person("Daan", "Wichmann", "Secret")
\end{minted}
\pause
\begin{minted}{python}
print(person.first_name)
>>> "Daan"
\end{minted}
\end{frame}

\begin{frame}[fragile]{Data Attributes}
Re-assigning an attribute:
\pause
\begin{minted}{python}
person = Person("Daan", "Wichmann", "Secret")
\end{minted}
\pause
\begin{minted}{python}
person.first_name = "Damai"
print(person.first_name)
>>> "Damai"
\end{minted}
\end{frame}

\begin{frame}[fragile]{Our beautiful class}
\begin{minted}{python}
class Person():
    
    # The constructor
    def __init__(self, first_name, last_name, age):
        self.first_name = first_name  # Attribute
        self.last_name = last_name  # Attribute
        self.age = age  # Attribute

    # Method
    def speak(self):
        print(f"Hey, my name is {self.first_name}")
\end{minted}
\end{frame}

\begin{frame}[fragile]{Calling our method}
\pause
\begin{minted}{python}
person = Person("Daan", "Wichmann", "Secret")
\end{minted}
\pause
\begin{minted}{python}
person.speak()
>>> Hey, my name is Daan
# Note that we don't fulfill the 'self' 
# parameter, but the interpreter
# does it automatically for us.
\end{minted}
\end{frame}

\begin{frame}{Why?}
Why do we want to use classes?
\begin{itemize}
    \item If we want to to represent something that is not already in Python
    \item To make collaboration easier
    \item If we want to 'group together' functionality
\end{itemize}
\end{frame}
